\section{Introduction}\label{sec:intro}
\sectionbrief
{noisy measurements $(t_i,x_i,y_i)$ with $y_i\approx u(t_i,x_i)$ of a field that solves an unknown linear constant-coefficient equation in one time and one space variable; the analytic theory (identification, factor decomposition, interaction term) is taken as given.}
{to turn exact theorems into a procedure whose input is imperfect data, and to say exactly where the exactness is lost.}
{this section frames the three tasks, shows the whole pipeline in one figure and lists what each stage relies on. Sections~\ref{sec:interp}--\ref{sec:loop} then describe the stages one by one.}

\subsection{The question}
Throughout, $u(t,x)$ is a field (a function of time $t$ and one space coordinate $x$) known on a rectangle $\Omega$ of the $(t,x)$-plane, and $P(\partial_t,\partial_x)u=0$ is the linear constant-coefficient equation it solves, with $P$ unknown; all the symbols are introduced in Section~\ref{sec:prelim}. The analytic work answers three questions for an exact solution $u$: which operator $P$ does $u$ determine (Q1), how does $u$ split into the components generated by the factors of $P$ (Q2), and what part of $u$ is due to no single term but to their interplay (Q3). Here the data are finitely many noisy samples, so the answers must change form:
\begin{enumerate}[label=\textbf{A\arabic*.},leftmargin=1.1cm,nosep]
\item \textbf{Which factors, with which coefficients, and how sure are we?} The output is a \emph{probability distribution over factor structures}, never a single ``identified'' law.
\item \textbf{How much of the signal does each factor carry?} Weights $w_j$ with credible intervals that widen when the components are hard to separate.
\item \textbf{Are the factors really separate or tangled?} The numbers $\gamma_N$, $c_F$ and the interaction ratio $\rho$, as distributions over interpolants.
\end{enumerate}
The setting is no longer purely analytic: part of the algorithm is statistical. The design rule is that the delicate analytic part stays deterministic and is called as a black box, and everything around it is Bayesian.

\subsection{The pipeline}
\begin{figure}[!htb]
\centering
\begin{tikzpicture}[node distance=7mm,
  det/.style={draw=cC,fill=cC!8,rounded corners=2pt,minimum height=15mm,text width=3.0cm,align=center,font=\small\sffamily},
  stat/.style={draw=cA,fill=cA!8,rounded corners=2pt,minimum height=15mm,text width=3.0cm,align=center,font=\small\sffamily},
  arr/.style={-{Stealth[length=2.4mm]},thick,cG}]
\node[stat] (d) {\textbf{0. Data}\\ $(t_i,x_i,y_i)$};
\node[stat,right=of d] (i) {\textbf{1. Interpolate}\\ GP posterior\\ $\{(\hat u_s,\pi_s)\}$};
\node[det,right=of i] (c) {\textbf{2. Spectral core}\\ modes $\to$ factors\\ per sample};
\node[stat,below=10mm of c] (a) {\textbf{3. Aggregate}\\ cloud $\mathcal W$, DP\\ posterior on $k$};
\node[stat,left=of a] (g) {\textbf{4. Group}\\ mode space,\\ rectangles};
\node[stat,left=of g] (b) {\textbf{5. Bayesian\\ output}\\ Friedrichs layer};
\node[stat,below=10mm of g] (l) {\textbf{6. Loop}\\ split, stop, regroup};
\draw[arr] (d)--(i); \draw[arr] (i)--(c); \draw[arr] (c)--(a); \draw[arr] (a)--(g); \draw[arr] (g)--(b);
\draw[arr] (g)--(l); \draw[arr,dashed] (l.west) -| ($(d.south)+(0,-4mm)$) -- (d.south);
\end{tikzpicture}
\caption{The pipeline. Only the green box is deterministic. The dashed arrow is the recursion: each group becomes a new data set.}\label{fig:pipeline}
\end{figure}

\begin{table}[!htb]
\centering\small
\begin{tabular}{@{}L{2.6cm}L{5.6cm}L{3.2cm}L{2.6cm}@{}}
\toprule
Stage & What it does & Relies on & Status\\\midrule
1 Interpolation & GP posterior on a grid from scattered data; hyperparameters marginalised & \thy{T11}, \thy{T12} & \impl{}; band coverage within $\pm5$ points in $11$ of $12$ cells\\
2 Spectral core & FFT in $x$, matrix pencil in $t$; structured factor fit, BIC over 85 structures; components, weights, $\gamma_N$; per-bin correlated noise & \thy{T1}, \thy{T2}, \thy{T5}, \thy{T15} & \impl{}\\
3 Aggregation & weighted cloud of factor coefficients; DP mixture; $k_{\mathrm{eff}}$ & \thy{T12} & \impl{}\\
4 Grouping & mode space, rectangles with quadtree & \thy{T1}, \thy{T5}, \thy{T8} & \impl{}\\
5 Bayesian output & posterior on components, $P(\text{split})$, intervals following $\theta_F$; raw-data evidence & \thy{T5}, \thy{T6} & \impl{}; coverage logged $0.85$--$1.00$, not certified\\
6 Loop & degree growth with a held-out guard, regroup at fixed $k$, recursion per factor, refinement on the observations & \thy{T15} & \impl{}; thresholds still hand defaults\\
6b Refinement and decision & raw-data evidence, support path, exact splits, goodness-of-fit downgrades, intervals (Sec.~\ref{sec:refine}) & \thy{T17}, T44--T103 & \impl{}; final benchmark $30/30$ right\\
7 Learned layers & seeds for the exponential fit, calibration, failure validator, remedy table (Phase~4; see the companion document) & \thy{T10} & evaluated; none ships except an exact memo; pipeline $=$ pipeline $+$ ML\\
\bottomrule
\end{tabular}
\caption{The stages, the theory each one rests on, and their present status.}\label{tab:stages}
\end{table}

\subsection{Notation at a glance}
Every symbol is introduced in Section~\ref{sec:prelim} or in the box that opens the section that needs it; this table only collects the main ones and says where each comes from.
\begin{table}[!htb]
\centering\small
\begin{tabular}{@{}L{2.9cm}L{7.0cm}L{3.0cm}@{}}
\toprule
Symbol & Meaning & Comes from\\\midrule
$u(t,x),\ \Omega$ & the field, and the rectangle of the $(t,x)$-plane where it is known & the data; Sec.~\ref{sec:prelim}\\
$P,\ Q_j$ & the unknown operator (as a polynomial in $\tau,\xi$) and its factors & the theory; Sec.~\ref{sec:prelim}\\
$\kappa,\ \lambda$ & wavenumber (spatial frequency, real) and rate of evolution in time (complex) of a mode $e^{\lambda t+i\kappa x}$ & Sec.~\ref{sec:prelim}\\
$m_j,\ dx_j,\ \mathcal M$ & degree in $\tau$ and bound on the degree in $\xi$ of factor $j$; a \emph{structure}, the list of these shapes (85 candidates) & Sec.~\ref{sec:prelim}; the algorithm\\
$\hat u_s,\ \pi_s,\ N_s$ & the $s$-th posterior field, its weight, the number of fields & Sec.~\ref{sec:interp}\\
$\sigma_\lambda$ & uncertainty of a mode's rate $\lambda$ (Cram\'er--Rao) & Sec.~\ref{sec:core}\\
$u_j,\ w_j$ & component of factor $j$; its weight $\|u_j\|/\sum_k\|u_k\|$ & Sec.~\ref{sec:core}\\
$\gamma_N,\ c_F$ & how well the components can be separated; $\gamma_2=1-c_F$ & theory; Sec.~\ref{sec:core}\\
$\rho$ & interaction ratio $\|\Eint\|/\|u(T)-u_0\|$ of an evolution factor & theory; Sec.~\ref{sec:core}\\
$\bs\beta^{(s)}_j$ & coefficient vector of factor $j$ in sample $s$ & Sec.~\ref{sec:agg}\\
$\mathcal W$ & the weighted cloud of all the $\bs\beta^{(s)}_j$ & Sec.~\ref{sec:agg}\\
$k_{\mathrm{eff}}$ & $\exp(-\sum_jW_j\log W_j)$, effective number of factors & Sec.~\ref{sec:agg}\\
\bottomrule
\end{tabular}
\caption{Notation guide.}
\end{table}

\subsection{Roadmap}
Section~\ref{sec:prelim} recalls why exponentials are the right ``basis''. Section~\ref{sec:interp} is stage 1, Section~\ref{sec:core} stage 2 (the longest), Section~\ref{sec:agg} stage 3, Section~\ref{sec:group} stage 4, Section~\ref{sec:bayes} stage 5 and Section~\ref{sec:loop} the loop. Section~\ref{sec:q} answers the two design questions put to the machine-learning side, Section~\ref{sec:train} describes the trainable version, Section~\ref{sec:example} follows one example from data to output, and Section~\ref{sec:limits} lists the limits.
